\mode*
\section{Conclusions}

% XXX Work in progress: conclusions to be written once the systematic
% review and the analysis rounds are complete. The restatements below keep
% the question--answer traceability in place from the start.

\rqmisconceptions*
So far: the seed round documents difficulty with composite commands
\autocite{Doane1992UnixPrompts}, the spatial nature of the file-system
skill \autocite{Bergman2008Navigation,Benn2015BrainFolders}, and a
research gap for the shell itself \autocite{Winder2024ShellTutor} and for
SSH (nothing found).

\rqaspects*
Thirteen candidate aspects, listed per objective in
\cref{sec:shell-aspects,sec:filesystem-aspects,sec:composition-aspects,sec:ssh-aspects}
and summarised in \cref{tab:aspects}.
Their evidential status is uneven and is stated with each: four are read
off difficulties documented in the literature, four off our own course
observations, and five are posited from the structure of the object of
learning; one of the thirteen (\emph{one tree, two views}) is the
file-system instance of another (\emph{the shell acts on the same system
as the GUI}).
Every aspect has an item in the pre/post-test of \cref{app:quiz} (two
have open twins as well: the first, since its closed form is asked only
at the end; \emph{a shell runs somewhere}, so that the account and the
choice can be compared), which is what will turn \enquote{candidate} into
\enquote{critical} or not.

\begin{table}
  \begin{sidecaption}{The candidate critical aspects, what each is read
    off, and the item of \cref{app:knowledge-items} that tests
    it.}[tab:aspects]
    \small
    \begin{tabular}{@{}>{\raggedright\arraybackslash}p{0.36\linewidth}l>{\raggedright\arraybackslash}p{0.34\linewidth}@{}}
      \toprule
      Aspect & Read off & Item \\
      \midrule
      \multicolumn{3}{@{}l}{\emph{The command model} (\cref{sec:shell-aspects})} \\
      A language, not a search box & literature & the command line reads a language; a space in a file name \\
      Textual feedback & literature & a message after a command \\
      Same system as the GUI & posited & the same files from two sides \\
      \addlinespace
      \multicolumn{3}{@{}l}{\emph{The file system} (\cref{sec:filesystem-aspects})} \\
      Everything has a place & literature & where a saved file is \\
      Where am I? & course & which directory the command uses \\
      One tree, two views & posited & the same files from two sides \\
      Absolute versus relative & course & two names for one file \\
      \addlinespace
      \multicolumn{3}{@{}l}{\emph{Composition} (\cref{sec:composition-aspects})} \\
      The pipe is not the file & literature, course & a file on the pipe; sending output into a pipe \\
      Uniform interfaces & posited & the same sort after different programs \\
      A script is nothing new & posited & running the same commands from a file \\
      \addlinespace
      \multicolumn{3}{@{}l}{\emph{The remote shell} (\cref{sec:ssh-aspects})} \\
      A shell runs somewhere & course & which computer answers (closed and open) \\
      Same shell, different computer & course & where exit leaves you \\
      Files live on machines & posited & where the remote file ends up \\
      \bottomrule
    \end{tabular}
  \end{sidecaption}
\end{table}

\rqpatterns*
Preliminary patterns accompany each aspect; they are to be realised as
concrete, captured command sequences and validated in teaching.

The natural design for that validation is a learning study
\autocite[Ch.~7]{NCOL} per objective, with the pre- and post-test of
\cref{app:quiz} as its instrument: the closed items test the aspects
argued for here, and the open accounts---the first phenomenographic data
on the shell since \textcite{DoyleLister2006Unix}---are where aspects we
did not think of will show themselves.
